% ============================================================
% 第 2 章 背景与相关工作
% ============================================================
\section{背景与相关工作}
\label{sec:related}

本章分两部分。前半部分（\S\ref{sec:rel-ashare}、\S\ref{sec:rel-sellside}）为不熟悉中国资本市场的读者补充领域背景：
A 股定期报告的结构与口径、卖方研报的体例与常用估值方法。这些细节不是装饰——\sysname{} 中相当一部分设计
（只取合并报表、营业收入优先于营业总收入、季报表头的期间识别、单季拆分规则、扣非口径、EBIT 定义）
都直接来自这些口径上的陷阱。后半部分（\S\ref{sec:rel-finqa}～\S\ref{sec:rel-harness}）按八个方向回顾相关工作，
每个方向说明本文与之的差异；\S\ref{sec:rel-compare} 用一张对比表总结定位。

% ============================================================
\subsection{领域背景：A 股定期报告}
\label{sec:rel-ashare}

\subsubsection{准则体系与年报结构}
A 股上市公司按《企业会计准则》\cite{mof2006cas} 编制财务报表，按中国证监会的信息披露内容与格式准则\cite{csrc2021annual}
编制年度报告。年报的章节顺序相对固定，表~\ref{tab:rel-annual} 列出与本文相关的部分。对系统设计而言，
最重要的事实是：\textbf{同一个指标会在年报中出现多次}，口径未必相同——“主要会计数据和财务指标”中的营业收入、
合并利润表中的营业收入、母公司利润表中的营业收入、附注中按产品拆分后的营业收入合计、管理层讨论与分析中的营业收入，
前两者应相等，第三者通常不同，第四者可能与前者有“其他业务收入”之差，第五者可能四舍五入为亿元。

\begin{table}[h]
\centering
\caption{A 股年度报告中与本文相关的章节（按通行顺序）。页码为一份约 250 页年报的典型位置，仅作示意。}
\label{tab:rel-annual}
\small
\begin{tabularx}{\linewidth}{@{}p{3.6cm}p{1.4cm}X@{}}
\toprule
章节 & 典型页码 & 对本文的用途 \\
\midrule
公司简介和主要财务指标 & 5～10 & “主要会计数据和财务指标”表（本期、上期、增减、上上期）；“分季度主要财务指标”表；非经常性损益项目及金额表 \\
管理层讨论与分析 & 10～60 & 主营业务、分产品 / 分地区 / 分渠道收入与毛利率、产销量、成本构成、费用、研发、现金流、资产负债变动、未来展望；经营数据（\file{extras.json}）的主要来源 \\
公司治理、环境与社会 & 60～100 & 一般不取数 \\
股份变动及股东情况 & 100～115 & 总股本、股本变动 \\
财务报告：审计报告 & 115～120 & 审计意见类型 \\
财务报告：合并报表 & 120～135 & 合并资产负债表、合并利润表、合并现金流量表、合并所有者权益变动表；真值表的主要来源 \\
财务报告：母公司报表 & 135～150 & 与合并报表科目同名，\textbf{必须排除} \\
财务报告：附注 & 150～250 & 会计政策、科目明细、现金流量表补充资料（折旧摊销）、分部信息、关联方；明细表与报表科目同名但口径不同 \\
\bottomrule
\end{tabularx}
\end{table}

\subsubsection{合并报表与母公司报表}
合并报表把母公司与其控制的子公司视为一个经济主体编制，母公司报表只反映母公司法人主体。
对控股型公司（例如经营在子公司、母公司主要持有股权的集团），母公司报表的营业收入可能远小于合并报表，
净利润则可能主要由子公司分红构成的投资收益决定。两套报表的科目名称几乎完全相同，在年报中先后出现，
一个只看“科目名 + 数字”的抽取器有相当概率取错。分析与估值只应使用合并报表，这是 \sysname{} 抽取环节的第一条硬约束：
模型抽取条目必须声明口径为合并（\code{scope=consolidated}），规则通道则按表格标题排除含“母公司”的表
（第 \ref{sec:extract} 节）。\code{locate} 命令会专门输出母公司报表所在页（\code{parent\_only\_pages}），提示模型避开。

\subsubsection{收入与利润的口径}
\paragraph{营业总收入与营业收入。}
一般企业的利润表第一行为“一、营业总收入”，其下“其中：营业收入”；对多数公司二者相等。
但若集团内有财务公司、保险或类金融业务，营业总收入还包括“利息收入”“已赚保费”“手续费及佣金收入”，
此时二者不等。研报中“营收”一般指营业收入。\sysname{} 的指标字典把二者都列为 \code{revenue} 的别名，
但规则通道在同一张表同时出现二者时只取“营业收入”行（\file{extract.py} 中的 \code{has\_op\_rev} 判断）。

\paragraph{净利润、归母净利润与扣非归母净利润。}
利润表中的“净利润”包含少数股东损益；“归属于母公司所有者的净利润”（年报关键指标表中写作“归属于上市公司股东的净利润”，
简称归母净利润）扣除少数股东损益，是每股收益与估值的基础；“归属于上市公司股东的扣除非经常性损益的净利润”
（简称扣非归母净利润或扣非净利润）再扣除非经常性损益。非经常性损益的认定依据证监会的解释性公告\cite{csrc2023nonrecurring}，
包括非流动资产处置损益、计入当期损益的政府补助（与正常经营密切相关、按标准持续享有的除外）、
金融资产公允价值变动与处置损益等。研报口语中“净利润”常被用来指归母净利润，这是纠错任务中“口径错误”的主要来源
（\S\ref{sec:intro-example} 的样例研报把含少数股东损益的净利润 6.46 亿元写成了归母净利润）。

\paragraph{EBIT 的口径。}
A 股利润表不列示“息税前利润”，且把投资收益、公允价值变动收益、资产处置收益、其他收益、信用与资产减值、
营业外收支都放在利润总额之前。估值所需的“经营性 EBIT”因此没有唯一定义。\sysname{} 采用
\begin{equation}
\text{EBIT} = \text{利润总额} + \text{财务费用} - \text{投资收益} - \text{公允价值变动收益},
\label{eq:rel-ebit}
\end{equation}
即剔除融资成本与投资类损益，保留资产处置、其他收益、减值与营业外收支（简化）。被剔除的投资类资产需要在股权价值桥中
作为非经营性资产加回，否则估值偏保守。这一口径在三表模型、DCF 与盈利预测摘要中保持一致（第 \ref{sec:problem}、\ref{sec:valuation} 节）。

\subsubsection{定期报告的期间口径}
A 股公司每年披露四份定期报告：一季报（1～3 月）、半年报（1～6 月）、三季报（1～9 月与 7～9 月）、年报（全年）。
其中\textbf{利润表与现金流量表只披露累计数}，单季数要靠相减得到；三季报的“主要会计数据”表通常同时有
“本报告期”（第三季度单季）与“年初至报告期末”（前三季度累计）两列，且列标题不写年份；
年报另有“分季度主要财务指标”表，直接列示四个单季的营业收入、归母净利润、扣非归母净利润与经营活动现金流量净额。
资产负债表是时点数，不能相减；每股收益与比率也不能相减。表~\ref{tab:rel-periods} 总结了这些期间之间的关系。

\begin{table}[h]
\centering
\caption{定期报告的期间与单季拆分关系。“发生额”指利润表、现金流量表与非经常性损益类指标；时点数、每股收益、比率不拆分。}
\label{tab:rel-periods}
\small
\begin{tabular}{@{}llll@{}}
\toprule
报告 & 披露期间（本文记号） & 表头常见写法 & 可推出的单季 \\
\midrule
一季报 & $Y$Q1 & 本报告期 & Q1 即单季 \\
半年报 & $Y$H1 & 本报告期 / 本期 & Q2 $=$ H1 $-$ Q1 \\
三季报 & $Y$Q3、$Y$Q1-3 & 本报告期 / 年初至报告期末 & Q3 $=$ Q1-3 $-$ H1（与列示的 Q3 互核） \\
年报 & $Y$；分季度表 Q1～Q4 & 本期 / 上期；第一～四季度 & Q4 $=$ 全年 $-$ Q1-3（与列示的 Q4 互核） \\
\bottomrule
\end{tabular}
\end{table}

这里有两个系统层面的陷阱。第一，“本报告期”在年报中指全年，在三季报中指第三季度单季，同一个词在不同文件中含义不同；
\sysname{} 的规则通道无法可靠区分，因此季报（通过 \code{--supplement} 传入）只接受模型抽取、不参与规则通道比对，
单季拆分由 \code{analyze()} 按表~\ref{tab:rel-periods} 的规则完成并登记公式。
第二，当单季值既有列示又能由累计数推出时，二者构成一项勾稽；但若某个累计数本身是由该单季推出的，
再拿来核对就是“自己和自己比”。\code{analyze()} 为每个单元格维护其依赖的“列示单元格集合”，据此排除这种退化核对
（第 \ref{sec:analysis} 节）。

\subsubsection{业绩预告、业绩快报与一致预期}
沪深交易所规则要求净利润为负、扭亏或大幅变动的公司在规定时限内发布业绩预告，通常给出归母净利润与扣非归母净利润的区间；
部分公司在年报前发布业绩快报。卖方“一致预期”是各家券商对同一指标预测值的汇总（均值或中位数），
由数据服务商发布。“实际 vs 一致预期”（超预期 / 低于预期，beat / miss）与“实际落在预告区间的位置”是年报点评的常规内容，
对应财报分析选题中的业绩分析部分。\sysname{} 不内置任何数据服务商接口：一致预期与业绩预告由人或研究助理会话以 CSV 形式提供，
后者需登记来源（第 \ref{sec:analysis}、\ref{sec:guards} 节）。

\subsubsection{版式层面的困难}
即便口径正确，从 PDF 中取数本身也有困难。表~\ref{tab:rel-layout} 列出我们在真实年报中遇到并处理的版式问题，
详细处理见第 \ref{sec:extract} 节。

\begin{table}[h]
\centering
\caption{真实 A 股年报 PDF 的版式问题与 \sysname{} 的处理。}
\label{tab:rel-layout}
\small
\begin{tabularx}{\linewidth}{@{}p{2.6cm}X X@{}}
\toprule
问题 & 表现 & 处理 \\
\midrule
单位切换 & 同一份年报中“元”“千元”“万元”“百万元”并存，单位写在表上方的“单位：”行 & 按页、按表识别单位提示；事实内部一律换算为元 \\
括号负数与占位符 & “(1,234.56)”表示负数；“—”“-”“不适用”表示无数 & 统一数字解析；占位符解析为缺失 \\
跨页续表 & 合并报表跨两三页，续页没有表头 & 续页沿用首页表头的期间列，并记录 \code{continued\_from} \\
稀疏网格 & pdfplumber 抽出的表格有大量空列、错位列 & 去除空列后再按表头对齐 \\
科目折行 & 长科目名（如政府补助的完整说明）被折成多行 & 合并折行后再匹配别名 \\
填列说明 & “投资收益（损失以“－”号填列）” & 别名匹配前去除填列说明 \\
同名明细表 & 附注中受限资产、费用明细、关联方往来等表的行名与报表科目同名 & 规则通道只从主要会计数据、分季度、非经常性损益、三大报表与补充资料取数 \\
追溯调整 & 表格同时列示“调整前 / 调整后” & 记为“追溯调整”问题，提示同比应使用调整后数据 \\
扫描页 & 无文本层 & 直接拒收，不做 OCR \\
\bottomrule
\end{tabularx}
\end{table}

% ============================================================
\subsection{领域背景：卖方研报体例与估值方法}
\label{sec:rel-sellside}

\subsubsection{研报体例与评级}
国内券商的深度报告在结构上高度趋同（\S\ref{sec:intro-scene}），首页的“投资评级”通常采用五档：买入、增持、中性、减持、卖出，
定义为未来 6～12 个月相对基准的预期超额收益区间，各家阈值略有差异。\sysname{} 的评级是一个确定性函数：
以基准情景 DCF 每股价值相对参考股价的上行空间 $u$ 为输入，$u\ge 20\%$ 买入、$u\ge 5\%$ 增持、$u\ge -5\%$ 中性、
$u\ge -20\%$ 减持、其余卖出（阈值可在假设文件中覆盖）。研报只引用程序给出的评级，不自行判断——这既避免了评级与估值的矛盾，
也把“观点”的来源明确限定在估值假设上。

\subsubsection{现金流折现：FCFF DCF}
企业自由现金流（free cash flow to the firm，FCFF）折现是卖方深度报告最常用的绝对估值方法\cite{damodaran2012investment,koller2020valuation}。
其核心关系为
\begin{align}
\text{FCFF}_t &= \text{EBIT}_t(1-\tau_t) + \text{D\&A}_t - \text{CapEx}_t - \Delta\text{NWC}_t, \label{eq:rel-fcff}\\
\text{EV} &= \sum_{t=1}^{N}\frac{\text{FCFF}_t}{(1+\text{WACC})^{t-0.5}} + \frac{\text{TV}_N}{(1+\text{WACC})^{N}}, \qquad
\text{TV}_N = \frac{\text{FCFF}_N(1+g)}{\text{WACC}-g}, \label{eq:rel-ev}
\end{align}
其中 $\tau_t$ 为税率，D\&A 为折旧摊销，NWC 为营运资本，$g$ 为永续增长率，预测期现金流按年中折现（折现期 $t-0.5$）、
终值按第 $N$ 年末折现；终值也可用退出倍数法 $\text{TV}_N=\text{EBITDA}_N\times$ 倍数。WACC 由资本资产定价模型给出的股权成本与税后债务成本按目标资本结构加权。
企业价值经股权价值桥（加现金、非经营性资产、联营投资，减有息负债、租赁负债、少数股东权益等）得到股权价值，再除以摊薄股本得到每股价值。
DCF 的结论对少数几个参数极其敏感（终值往往占企业价值的一半以上），因此卖方报告通常给出 WACC × 永续增长率的敏感性矩阵与多情景结果。

\subsubsection{三表联动}
DCF 中的 EBIT、折旧摊销、资本开支与营运资本变动本应来自一个完整的财务预测模型：利润表、资产负债表、现金流量表三表联动，
资产负债表每期平衡、现金流量表的期末现金与资产负债表现金一致、留存收益按净利润与分红滚动、利息费用与债务余额相互依赖（通常需要迭代求解）。
很多简化的 DCF 实现直接假设“EBIT 利润率”“资本开支占收入比”等向量，跳过三表模型；这使 DCF 与盈利预测表可能互不一致
（例如盈利预测表中的净利润与 DCF 隐含的 NOPAT 对不上）。\sysname{} 让 DCF 的全部经营向量逐期取自三表模型的同一情景，
并把“DCF 与三表同口径”作为评分标准的硬闸门（第 \ref{sec:valuation}、\ref{sec:eval} 节）。

\subsubsection{可比公司估值}
可比公司法（comparable companies，comps）用同业的估值倍数（P/E、EV/EBITDA、EV/Sales、P/B 等）推算目标公司价值\cite{liu2002multiples}。
它的难点在于可比公司的选取、倍数的口径配对（EV 类倍数配 EV 类分母）、亏损或负值分母的处理（记为 NM，不参与统计）
以及离群值处理（常用四分位距规则）。在卖方报告中，可比公司法通常作为 DCF 的交叉验证，并解释目标公司相对同业溢价或折价的原因。

\subsubsection{反向 DCF}
反向 DCF（reverse DCF）由 Rappaport 与 Mauboussin 的“预期投资”框架推广\cite{rappaport2001expectations}：与其从假设推出价值，
不如固定价格、反推“当前股价隐含了什么样的增长与盈利能力”，再判断这些隐含预期是否现实。它直接服务于研报中的“预期差”论证——
如果股价隐含的收入复合增速远高于历史与行业水平，看多就需要更强的理由。\sysname{} 的反向 DCF 固定其余假设为基准情景，
每次只动一个变量（收入增速、EBIT 利润率、WACC、永续增长率或退出倍数），用网格扫描找变号区间再二分，
求出“每股价值恰好等于参考股价”时该变量的取值，并按历史区间给出可行性分档（第 \ref{sec:valuation} 节）。

% ============================================================
\subsection{金融数值推理与财报问答}
\label{sec:rel-finqa}

财报问答是最早系统研究“LLM 在财务数字上推理”的方向。FinQA\cite{chen2021finqa} 从标普 500 公司的财报中构造了约八千个需要多步数值推理的问答对，
并为每题标注了由加减乘除、比较等运算组成的推理程序；ConvFinQA\cite{chen2022convfinqa} 把它扩展为多轮对话，考察推理链在对话中的延续；
TAT-QA\cite{zhu2021tatqa} 关注表格与文本混合的上下文，答案需要从表格中取数、从文本中取条件；
FinanceBench\cite{islam2023financebench} 以真实公开文件为语料构造开卷问答，报告了带检索的强模型在其样本上仍有大量问题答错或拒答；
DocFinQA\cite{reddy2024docfinqa} 把 FinQA 的上下文从摘录扩展到完整文件，长度达到十万词量级；
BizBench\cite{koncel2024bizbench} 则把商业与金融定量推理整理为包含程序合成的基准。PIXIU\cite{xie2023pixiu} 与 FinBen\cite{xie2024finben}
把这些任务连同情感分析、信息抽取、预测等任务汇总成金融大模型的综合评测。

\paragraph{与本文的关系。}
这些工作确立了两个本文赖以出发的观察：模型在财务数字上的多步推理不可靠；把推理写成可执行程序能显著提升正确率。
但它们的设定与研报生产有三点根本不同。(1)~\textbf{评价单位不同}：财报问答的单位是“一个问题一个答案”，
研报的单位是“一篇数万字、含数百个相互依赖数字的文稿”，错误可以出现在任何一个数字上，也可以出现在数字之间的关系上（前后矛盾、构成描述）。
(2)~\textbf{输出约束不同}：问答只要求答案正确，研报还要求每个数字可追溯到来源、区分事实推论观点、改写后不变。
(3)~\textbf{语言与口径不同}：上述基准以美股英文财报为主，A 股中文年报的口径陷阱（合并与母公司、营业总收入、扣非、季报表头）
与版式问题（单位切换、跨页续表）不在其覆盖范围内。本文的数据提取、财报分析与纠错核查可以看作“面向 A 股的、以证据定位为必要条件的财报问答与核查”，
但我们不构造问答基准，而是在系统中把“答对”强化为“能证明答对”。

% ============================================================
\subsection{程序辅助推理与工具使用}
\label{sec:rel-tools}

思维链（chain-of-thought）提示\cite{wei2022cot} 与自洽性解码\cite{wang2023selfconsistency} 提升了模型的多步推理，但算术仍由模型完成。
PAL\cite{gao2023pal} 与 Program of Thoughts（PoT）\cite{chen2023pot} 让模型生成程序、由解释器执行，把“推理”与“计算”分开，
在数学与金融数值推理上大幅优于纯文本推理。Toolformer\cite{schick2023toolformer} 让模型自监督地学会何时调用计算器、检索等工具；
ReAct\cite{yao2023react} 把推理与行动交错，形成“思考—调用工具—观察”的循环，成为后来大多数智能体框架的基本模式。

\paragraph{与本文的关系。}
\sysname{} 继承了 PAL / PoT 的核心洞见——计算应交给确定性执行器——但在两个方向上走得更远。
\textbf{其一，计算不是模型可选调用的工具，而是一条固定的流水线。} 在 PAL 中，模型决定写什么程序、程序算什么；
在 \sysname{} 中，算什么是预先确定的（三条流水线的步骤由 \code{Pipeline.step} 注册），模型只提供输入（抄录的数字与估值假设），
不编写计算代码。这消除了“程序写错”这一新的失败面，也让计算结果在不同运行之间可比、可复现。
\textbf{其二，约束模型的输出而不仅是模型的计算。} 工具使用框架只保证“调用了工具的那部分”是对的，
对模型在不调用工具时直接写下的数字没有约束；\sysname{} 的数字护栏对文稿中的每个数字做出处核对，
不论它是怎么产生的。换言之，PAL 回答的是“模型如何正确地算”，本文回答的是“如何保证模型写下的数都是算出来或抄对的”。

\begin{tradeoff}[：让模型写计算代码 vs 固定计算流水线]
让分析师会话直接写 Python 计算增速与估值，灵活性最高：遇到新的派生指标无须改引擎。我们没有这样做，原因有三：
(1)~模型写的计算代码本身需要审计，而审计一段每次都不同的代码比审计一个固定引擎困难得多；
(2)~评审复现要求“用保存的输入文件逐位复算”，若计算代码由模型生成，复现就依赖模型输出；
(3)~估值口径（EBIT 定义、年中折现、股权价值桥的项目）需要在全部公司之间一致，固定引擎天然保证这一点。
作为折中，\file{extras.json} 允许模型声明“我需要某两个已核验数字的增速或差值”（\code{ratios} 中的 \code{op: growth / diff}），
由程序计算并登记出处——模型可以决定“算什么”，但不能决定“怎么算”和“结果是多少”。
\end{tradeoff}

% ============================================================
\subsection{金融大模型与金融智能体}
\label{sec:rel-finagents}

\paragraph{金融大模型。}
BloombergGPT\cite{wu2023bloomberggpt} 以大规模金融语料与通用语料混合训练了一个约五百亿参数的模型，在金融 NLP 任务上优于同规模通用模型；
FinGPT\cite{yang2023fingpt} 走开源、数据中心的路线，强调低成本的轻量适配与持续更新的金融数据管道。
这一方向的目标是让模型本身“更懂金融”。本文不训练模型；我们的经验是，对研报生产而言，现成通用模型的金融知识已经足够，
瓶颈在于数字的可靠性与可追溯性，而这不是通过领域预训练能解决的问题——一个更懂金融的模型同样可能取到母公司报表的数字。

\paragraph{金融智能体。}
FinMem\cite{yu2023finmem} 为交易智能体设计了分层记忆与角色设定；Zhang 等人的 FinAgent\cite{zhang2024finagent} 是一个面向交易的多模态基础智能体，
结合工具增强与多样化的市场信息\footnote{与本系统的代码仓库名 \code{finagent} 同名，二者无关。}；
TradingAgents\cite{xiao2024tradingagents} 模拟交易公司的组织结构，设置基本面、情绪、新闻、技术分析师与多空研究员、交易员、风控等角色，
通过辩论形成交易决策；FinRobot\cite{yang2024finrobot} 是一个开源的金融智能体平台，提供分层的智能体与工具组织，
其示例包括基于公司年报生成个股分析报告。

\paragraph{与本文的关系。}
这些系统大多以\textbf{交易决策}为目标，评价指标是回测收益、夏普比率等；本文以\textbf{研究产出}为目标，
评价指标是数字闭合、分析质量与表达合规，不做任何交易或收益声明。即便是与本文最接近的研报生成类示例，
其关注点也在于“能否生成一篇结构完整的报告”，对“报告中每个数字从哪里来、错了能否被发现”缺少系统性的约束——
数字通常由模型从检索结果或工具返回中直接转写进文本，既没有逐字回查，也没有对最终文稿的出处核对。
此外，现有金融智能体中的多角色通常在同一进程、同一对话历史中协作（辩论、讨论），而本文刻意让角色之间\textbf{不共享上下文}，
以获得独立复核的效果（\S\ref{sec:rel-multiagent}）。

% ============================================================
\subsection{多智能体框架}
\label{sec:rel-multiagent}

AutoGen\cite{wu2023autogen} 提供可对话智能体的抽象，智能体之间通过消息对话协作，并可接入工具与人类；
CAMEL\cite{li2023camel} 通过角色扮演与“起始提示”让两个智能体自主协作完成任务；
MetaGPT\cite{hong2024metagpt} 把软件公司的标准作业流程（SOP）编码进多智能体系统，各角色产出结构化文档（需求、设计、接口）而不是自由对话，
以减少级联幻觉；ChatDev\cite{qian2024chatdev} 以“聊天链”组织设计、编码、测试等阶段；
Generative Agents\cite{park2023generative} 研究了带记忆与反思的智能体社会模拟；
多智能体辩论\cite{du2024debate} 让多个模型实例相互质疑以提升事实性与推理。

\paragraph{与本文的关系。}
\sysname{} 与 MetaGPT 在精神上最接近：都把一个人类组织的生产流程映射为多个角色，都以结构化产物而不是对话作为交接物。
差异在于三点。
\textbf{(1) 编排器不是模型。} 上述框架中，任务的推进、角色的切换、何时结束通常由某个模型（或模型之间的对话）决定；
\sysname{} 的编排器是确定性 Python 状态机，起哪个会话、闸门是否通过、问题单路由给谁、何时停止都写成代码，行为可预期、可单元测试。
\textbf{(2) 交接物受程序闸门约束。} MetaGPT 的结构化文档提高了交接质量，但文档内容是否正确仍由下游模型判断；
\sysname{} 的每个交接文件都有程序判定的闸门（数字逐字回查、来源摘句含数、护栏通过、数字锁通过、页数与图表数下限），
不通过就不会启动下一个角色。
\textbf{(3) 上下文隔离是设计目标而不是实现细节。} 多数框架让角色共享对话历史或共享记忆，以便协作；
我们观察到共享上下文会让复核者继承写作者的盲点、让评审者因看到过程而宽容，因此每个角色是一次独立的无头会话，
连工作区都与仓库物理隔离，盲评更在第三个目录中进行。多智能体辩论依赖“不同实例给出不同观点”，
本文则依赖“不同角色看到不同信息”——复核编辑只看到文稿与产物，看不到写作过程；盲评只看到交付物，看不到任何过程。

% ============================================================
\subsection{自我改进与反思}
\label{sec:rel-refine}

Self-Refine\cite{madaan2023selfrefine} 让同一模型对自己的输出给出反馈并迭代改写；Reflexion\cite{shinn2023reflexion} 让智能体把失败经验写成语言反思存入记忆，
在后续尝试中利用；Chain-of-Verification\cite{dhuliawala2024cove} 让模型先起草、再生成验证问题、独立回答后修正初稿。
另一方面，Huang 等人\cite{huang2024cannot} 指出，在没有外部反馈的情况下，模型的“自我纠正”往往不能提升推理正确率，甚至会把对的改错。

\paragraph{与本文的关系。}
\sysname{} 的 \code{refine} 闭环在形式上是一个“生成—评价—修订”的迭代，但它在三个方面区别于自我改进：
(1)~\textbf{反馈来自外部且部分是程序}——L0/L1 由程序判定，L2/L3 由一个看不到写作过程的不同会话判定，
这与 Huang 等人的结论一致：有效的改进需要外部信号；
(2)~\textbf{修订受数字锁约束}——Self-Refine 式的改写可以任意修改输出，而本文中复核与修订会话的改写必须保持数字不增不改，
需要改数字的问题被路由回分析师重跑数据阶段，由引擎重新计算；
(3)~\textbf{有停止与回滚规则}——迭代满三轮、连续两轮提升不足 2 分或无可回传问题时停止，结束时回滚到得分最高的一轮，
避免“越改越差”。Reflexion 的跨回合记忆在本文中没有对应物：每轮修订会话都是新会话，只读脱敏的问题单，
不积累对评分标准的“经验”，这是有意为之，以减少对评分标准的过拟合。

% ============================================================
\subsection{幻觉与可追溯生成}
\label{sec:rel-attribution}

幻觉（hallucination）是生成模型的系统性问题\cite{ji2023survey}。检索增强生成（RAG）\cite{lewis2020rag} 把外部文档作为条件，
Self-RAG\cite{asai2024selfrag} 让模型自行决定何时检索并对生成内容自我评判。在“可追溯”方向，
ALCE\cite{gao2023alce} 建立了带引用文本生成的基准，用引用召回与引用精确率衡量引用是否支撑陈述；
Attributed QA\cite{bohnet2022attributed} 与 AIS 框架\cite{rashkin2023attribution} 形式化了“可归因于来源”的判定；
Liu 等人\cite{liu2023verifiability} 评测了生成式搜索引擎，发现相当比例的陈述没有被其引用完全支撑；
FActScore\cite{min2023factscore} 把长文拆成原子事实逐条核验；RARR\cite{gao2023rarr} 在生成后检索证据并修订不被支撑的内容。

\paragraph{与本文的关系。}
本文的“数字可追溯”可以看作可归因生成在数值领域的一个强化版本，有三点不同。
\textbf{(1) 判定是确定性的。} ALCE、AIS、FActScore 判断“引用是否支撑陈述”依赖自然语言推理模型或人工；
\sysname{} 只对数字做判定，判定规则是字符串与数值比较：抄录的数字必须按数字边界出现在所写页码的原文中，
外部登记的数字必须原样出现在摘句中，文稿中的数字必须按展示精度等于某个产物值。这种判定无法覆盖非数值陈述，
但在其覆盖范围内没有判定误差。
\textbf{(2) 来源有层级。} 一般的引用生成只有“文档—陈述”一层；本文的数字有“原文—事实—真值表—派生—产物—文稿”多层，
每层都可以反向追溯，派生数的出处是公式与输入单元格而不是某段原文。
\textbf{(3) 出处类型被显式区分。} 本文按出处把文稿中的每个数字归为事实（报告披露、外部来源）、推论（程序计算）或观点（模型预测、估值假设），
这对应竞赛与行业对“区分事实、推论与观点”的要求，是一般可归因生成框架没有的维度。
另一方面，我们的出处核对只证明“存在一个匹配的值”，不证明“该值配给了正确的指标”——后者由指标—期间绑定核对部分覆盖，
其余依赖复核与盲评。这一局限在第 \ref{sec:problem} 节的威胁模型中明确列出。

% ============================================================
\subsection{LLM 评审与基于评分细则的评测}
\label{sec:rel-judge}

用强模型作为评审（LLM-as-a-judge）已成为开放式生成评测的主流做法。Zheng 等人\cite{zheng2023judging} 以 MT-Bench 与 Chatbot Arena 系统研究了
LLM 评审与人类偏好的一致性，并指出位置偏差、冗长偏好与自我增强偏差；Wang 等人\cite{wang2024fair} 进一步展示了仅调换候选顺序即可改变评审结论；
Panickssery 等人\cite{panickssery2024self} 发现评审模型能识别并偏好自己的生成。G-Eval\cite{liu2023geval} 用思维链与表单填写提升评审与人类的一致性；
Prometheus\cite{kim2024prometheus} 训练了以细粒度评分细则（rubric）为条件的开源评审模型；HealthBench\cite{arora2025healthbench}
以大量医生撰写的、针对具体对话的评分条目评测医疗对话，把“开放式质量”拆成可逐条判定的二元条目，再由程序加权计分。

\paragraph{与本文的关系。}
本文的评测体系在结构上借鉴了 HealthBench 一类的“逐条二元判定 + 程序计分”，在此基础上做了四点调整。
\textbf{(1) 程序判与模型判分层。} 能由程序确定判定的条目（数据闭合、数字前后矛盾）一律由程序判，
模型只判程序判不了的分析质量与表达；v2 评分标准把“数字前后矛盾”从模型判改为程序判，正是因为观察到评审对该条目的判断随措辞波动。
\textbf{(2) 评分标准对生成方不可见。} 生成方的工作区不含评分标准与质检代码，问题单只含理由与证据，不含条目原文与分值，
以减轻 Goodhart 效应\cite{manheim2018goodhart}。
\textbf{(3) 评审看不到过程与身份。} 盲评目录只含交付物，评审不知道写作过程与生成模型名，以减轻自我偏好与“过程宽容”。
\textbf{(4) 评审本身被校准。} 我们不假定评审可靠，而是用三组证据度量它：同稿重复评审的逐条一致率与总分极差；
程序制造已知缺陷（劣化算子）后的判出率与对非目标条目的连带改判；人工复核评审判出问题的精确率。
这些度量的结果并不完美（\S\ref{sec:intro-results}），但它们使“评审有多可靠”成为一个可报告的数字，而不是默认假设。

% ============================================================
\subsection{编码智能体 harness}
\label{sec:rel-harness}

近两年，面向软件工程的编码智能体发展迅速。SWE-bench\cite{jimenez2024swebench} 以真实 GitHub 问题评测模型修复代码的能力；
SWE-agent\cite{yang2024sweagent} 提出“智能体—计算机接口”（agent-computer interface），说明为模型设计合适的命令与反馈格式能显著提升任务完成率；
OpenHands\cite{wang2025openhands} 提供了开放的通用编码智能体平台。Claude Code\cite{anthropic2025claudecode} 是 Anthropic 发布的命令行编码智能体，
支持读写文件、执行命令、联网检索、技能（skill）文档与项目级 \file{CLAUDE.md} 指令，并提供无头模式（\code{claude -p}）
以 JSON 事件流输出完整的工具调用记录，以及按工具与路径的权限白名单 / 黑名单。

\paragraph{与本文的关系。}
本文把一个通用编码智能体用作投研多智能体系统的 harness，而不是从头编写智能体循环。这一选择的理由是：
(1)~研报生产中模型要做的事——读 PDF、读写 JSON 与 Markdown、运行命令行校验、根据报错修改文件、联网检索、写 Python 排版脚本——
恰好是编码智能体最擅长的工作模式；(2)~SWE-agent 的经验表明，接口设计比智能体循环本身更重要，
我们因此把精力放在“给模型什么命令、命令返回什么反馈”上：每个引擎命令都返回结构化的拒收理由与下一步建议，
\code{status} 命令直接告诉会话“下一道闸门是什么、差什么”；(3)~无头模式的事件流天然满足竞赛“全程记录工具调用、可追溯”的要求，
编排器把它原样落盘到工作区 \file{.pipeline/agents/}；(4)~权限系统让我们能在操作系统层面之上再加一层约束：
禁止读取仓库目录、禁止修改引擎与技能。
与 SWE-bench 一类工作不同，我们不评测 harness 本身的能力，而是把它当作一个可替换的执行器——编排器只依赖“启动一个会话、
传入提示与工具白名单、读回事件流”这一最小接口。

% ============================================================
\subsection{对比与定位}
\label{sec:rel-compare}

表~\ref{tab:rel-compare} 把代表性工作与本文在八个能力维度上做对比。需要说明：这些工作的目标各不相同，
多数并不以表中维度为设计目标，表格只用于说明本文的定位，不构成优劣评价；“$\circ$”表示部分具备或依配置而定。

\begin{table}[t]
\centering
\caption{与代表性工作的能力对比。A：计算由固定的确定性流水线完成（而非模型心算或模型临时编写的代码）；
B：模型给出的数字逐字回查原文页面；C：外部信息登记来源且正文强制引用；D：改写受数字锁约束；
E：角色之间上下文隔离；F：评分标准对生成方不可见；G：评审本身被校准；H：面向 A 股年报到研报的端到端产出。
\cmark{} 具备，$\circ$ 部分具备，\xmark{} 不具备或不涉及。}
\label{tab:rel-compare}
\small
\setlength{\tabcolsep}{4pt}
\begin{tabular}{@{}lcccccccc@{}}
\toprule
工作 & A & B & C & D & E & F & G & H \\
\midrule
FinQA / TAT-QA / FinanceBench\cite{chen2021finqa,zhu2021tatqa,islam2023financebench} & \xmark & \xmark & \xmark & \xmark & \xmark & \xmark & \xmark & \xmark \\
PAL / PoT\cite{gao2023pal,chen2023pot} & $\circ$ & \xmark & \xmark & \xmark & \xmark & \xmark & \xmark & \xmark \\
Toolformer / ReAct\cite{schick2023toolformer,yao2023react} & $\circ$ & \xmark & $\circ$ & \xmark & \xmark & \xmark & \xmark & \xmark \\
BloombergGPT / FinGPT\cite{wu2023bloomberggpt,yang2023fingpt} & \xmark & \xmark & \xmark & \xmark & \xmark & \xmark & \xmark & \xmark \\
FinMem / FinAgent（交易）\cite{yu2023finmem,zhang2024finagent} & $\circ$ & \xmark & \xmark & \xmark & \xmark & \xmark & \xmark & \xmark \\
TradingAgents\cite{xiao2024tradingagents} & $\circ$ & \xmark & \xmark & \xmark & \xmark & \xmark & \xmark & \xmark \\
FinRobot\cite{yang2024finrobot} & $\circ$ & \xmark & $\circ$ & \xmark & \xmark & \xmark & \xmark & $\circ$ \\
AutoGen / CAMEL\cite{wu2023autogen,li2023camel} & \xmark & \xmark & \xmark & \xmark & \xmark & \xmark & \xmark & \xmark \\
MetaGPT / ChatDev\cite{hong2024metagpt,qian2024chatdev} & \xmark & \xmark & \xmark & \xmark & $\circ$ & \xmark & \xmark & \xmark \\
Self-Refine / Reflexion\cite{madaan2023selfrefine,shinn2023reflexion} & \xmark & \xmark & \xmark & \xmark & \xmark & \xmark & \xmark & \xmark \\
ALCE / Self-RAG / RARR\cite{gao2023alce,asai2024selfrag,gao2023rarr} & \xmark & $\circ$ & \cmark & \xmark & \xmark & \xmark & \xmark & \xmark \\
MT-Bench / Prometheus / HealthBench\cite{zheng2023judging,kim2024prometheus,arora2025healthbench} & \xmark & \xmark & \xmark & \xmark & \xmark & $\circ$ & $\circ$ & \xmark \\
SWE-agent / OpenHands / Claude Code\cite{yang2024sweagent,wang2025openhands,anthropic2025claudecode} & \xmark & \xmark & \xmark & \xmark & $\circ$ & \xmark & \xmark & \xmark \\
\midrule
\textbf{\sysname{}（本文）} & \cmark & \cmark & \cmark & \cmark & \cmark & \cmark & \cmark & \cmark \\
\bottomrule
\end{tabular}
\end{table}

从表中可以看出本文的定位：本文没有在任何单一维度上提出全新的技术原语——确定性计算、引用核验、多角色、LLM 评审都已有大量前人工作；
本文的贡献在于\textbf{把这些原语以“不变量 + 强制机制”的方式组合进同一个系统}，使它们在研报生产这一长链条任务中同时成立，
并且对每一条不变量都给出执行它的具体代码位置与测试。这种组合并不平凡：例如数字锁（D）只有在数字护栏（A、B、C）已经保证初稿数字有出处的前提下才有意义；
评分标准不可见（F）只有在工作区物理隔离、问题单脱敏的前提下才能成立；评审校准（G）需要能够程序化制造已知缺陷的劣化算子，
而这又依赖研报纠错引擎（纠错核查子选题）对文稿结构的解析能力。各机制之间的依赖关系在第 \ref{sec:principles} 节展开。

% ============================================================
\subsection{小结}
\label{sec:rel-summary}

本章补充了理解后文所需的两类背景。在领域侧，A 股定期报告的口径（合并与母公司、营业总收入与营业收入、归母与扣非、
季报表头的期间含义、单季拆分）与版式（单位切换、跨页续表、同名明细表）共同决定了“从年报取对一个数”远比看上去困难；
卖方研报的估值方法（三表联动、FCFF DCF、可比公司、反向 DCF）对少数参数高度敏感，且各部分之间必须口径一致。
在方法侧，财报问答与程序辅助推理证明了“计算应交给执行器”，可归因生成证明了“陈述应有可核验的来源”，
多智能体框架证明了“按组织流程分工”有效，LLM 评审研究提醒我们“评审本身有偏差、需要度量”。
本文在这些工作之上，针对研报生产这一长链条、高数字密度、需要可追溯与可复现的任务，把上述思想组合为一组可强制执行的系统不变量。
下一章给出这些不变量的形式化定义。
